% Values transcribed from the P6-accepted expected_fragments.tsv.
Table~\ref{tab:screening} and Figure~\ref{fig:bands} summarize screening. Flanking PCR distinguishes insert-bearing candidates from empty vector. A negative orientation assay is read together with an amplification-positive control. If all orientation assays are negative in the first round, flanking-positive candidates go to sequencing. Reverse TFEB is a hypothetical challenge; incompatible-end directional ligation does not favor it.

\begin{table}[htbp]
\centering\small
\caption{Predicted screening products. All sizes are nominal bp; ``none'' means no specific product is predicted. R and S are hypothetical challenge sequences.}
\label{tab:screening}
\begin{tabularx}{\textwidth}{@{}L{3.2cm}YYY@{}}
\toprule
Sample & Flanking PCR (S-F1/S-R1) & Orientation PCR (S-F2/S-R1) & HindIII/BamHI digest \\
\midrule
Intended construct (C) & 2,605 & 2,093 & 4,695 + 1,439 \\
Empty vector (E) & 1,204 & none & 4,695 + 38 \\
Reverse TFEB (R) & 2,605 & none & 4,695 + 1,439 \\
Retained TFEB stop (S) & 2,608 & 2,096 & 4,695 + 1,442 \\
\bottomrule
\end{tabularx}
\end{table}

Flanking PCR and HindIII/BamHI digestion do not separate reversed TFEB from the intended construct, and the three-base retained-stop difference falls below gel resolution. Acceptance therefore rests on the expected digest plus actual, unambiguous, concordant bidirectional sequence across 602--2799 (2,198 nt), including both junctions, the complete fusion ORF and terminal stop. Ten primers predict usable coverage in both directions at every target base. Low-quality calls, mixed peaks, gaps and unexpected target changes are resolved under the Sanger acceptance procedure before cell assays. Figure~\ref{fig:workflow} summarizes the planned workflow and controls.
